\subsubsection*{Evaluation}

\subsubsection*{Evaluation Setup}
Model performance was evaluated on a held-out test set covering the period January 2023 through December 2025. A strict temporal holdout was used, ensuring that all test observations occur after the training period.

Performance was assessed using multiple complementary metrics, including coefficient of determination ($R^2$), mean absolute error (MAE), root mean squared error (RMSE), unbiased RMSE (ubRMSE), bias, and quantile-based error statistics such as the 90th percentile absolute error.

\subsubsection*{Overall Performance}

\begin{table}[H]
\centering
\begin{tabular}{lc}
\toprule
\textbf{Metric} & \textbf{Value} \\
\midrule
$R^2$ & 0.822 \\
MAE & 0.0283 \\
RMSE & 0.0397 \\
ubRMSE & 0.0396 \\
Bias & -0.0032 \\
P90 $|$error$|$ & 0.0615 \\
\bottomrule
\end{tabular}
\caption{Performance of the weighted base model on the test set (2023–2025).}
\end{table}

These results indicate that the model captures a substantial portion of the variability in soil moisture while maintaining low absolute and squared error.

\subsubsection*{Effect of Temporal Weighting}
To assess the impact of temporal weighting, the weighted model was compared to an unweighted baseline trained with identical hyperparameters.

\begin{table}[H]
\centering
\begin{tabular}{lcc}
\toprule
\textbf{Metric} & \textbf{Unweighted} & \textbf{Weighted} \\
\midrule
$R^2$ & 0.814 & 0.822 \\
MAE & 0.0304 & 0.0283 \\
RMSE & 0.0406 & 0.0397 \\
ubRMSE & 0.0405 & 0.0396 \\
Bias & -0.0028 & -0.0032 \\
P90 $|$error$|$ & 0.0634 & 0.0615 \\
\bottomrule
\end{tabular}
\caption{Comparison of unweighted and temporally weighted models.}
\end{table}

Temporal weighting leads to consistent improvements across most metrics, including higher $R^2$ and lower error values. This suggests that emphasizing more recent observations improves generalization to future data, likely due to nonstationarity in environmental conditions.

\subsubsection*{Residual Analysis and Heteroscedasticity}
Residual diagnostics revealed that model errors are not uniformly distributed across the target range. Both Breusch–Pagan and White tests indicate statistically significant heteroscedasticity (p-values $\ll 0.001$), confirming that error variance depends on soil moisture magnitude.

Visual inspection further shows structured residual behavior, with increased variance in specific regions of the target space. This suggests that a single global model may not fully capture the variability of the underlying physical processes.

\subsubsection*{Regime-Based Performance}
To further investigate error structure, the test data were partitioned into quantile-based bins according to the true soil moisture values, and model performance was evaluated within each bin.

The results show substantial variation across regimes. In particular, mid-range moisture conditions exhibit significantly worse fit compared to very dry conditions, with some bins producing strongly negative $R^2$ values.

This pattern indicates that the modeling task is inherently regime-dependent, with different physical processes dominating across moisture ranges. These findings directly motivate the exploration of regime-specific modeling approaches in the subsequent section.
